\documentclass[11pt]{article}

\usepackage[T1]{fontenc}
\usepackage{amsmath,amssymb,amsthm}
\usepackage{hyperref}

\title{Asymptotics for Guttman's Multiple-Group CFA Estimator}
\author{}
\date{\today}

\newcommand{\T}{^{\mathsf T}}
\newcommand{\E}{\mathbb{E}}
\newcommand{\Var}{\operatorname{Var}}
\newcommand{\Cov}{\operatorname{Cov}}
\newcommand{\tr}{\operatorname{tr}}
\newcommand{\diag}{\operatorname{diag}}
\newcommand{\vech}{\operatorname{vech}}
\newcommand{\vecop}{\operatorname{vec}}
\newcommand{\IF}{\operatorname{IF}}
\newcommand{\argmin}{\operatorname*{arg\,min}}

\theoremstyle{plain}
\newtheorem{proposition}{Proposition}
\theoremstyle{remark}
\newtheorem*{remark}{Remark}

\begin{document}
\maketitle

\section*{Purpose}

Dhaene and Rosseel \cite{dhaene2024} report that Guttman's multiple-group
method \cite{guttman1952} is a surprisingly strong non-iterative CFA estimator
in their simulations. This note records the asymptotic objects needed to
evaluate that claim analytically against normal-theory ML and the
moment-quadratic estimators already implemented in \texttt{magmaan}.

The phrase ``multiple-group method'' below is Guttman's historical name for
forming factor score groups from indicators. It is not multi-group SEM. Actual
multi-group SEM enters only by stacking independent covariance blocks.

\section{Sample Moment Setup}

Let \(y_i\in\mathbb{R}^p\) have mean \(\mu_0\), covariance \(\Sigma_0\), and
finite fourth moments. Write
\[
  \hat\sigma=\vech(\hat\Sigma),\qquad
  \sigma_0=\vech(\Sigma_0),
\]
with the usual covariance influence rows
\[
  z_i=\vech\{(y_i-\mu_0)(y_i-\mu_0)\T-\Sigma_0\}.
\]
Then
\[
  \sqrt n(\hat\sigma-\sigma_0)
    = n^{-1/2}\sum_{i=1}^n z_i+o_p(1),
  \qquad
  \Gamma=\Var(z_i).
\]
Under multivariate normality, \(\Gamma\) reduces to the normal-theory
covariance of \(\vech(\hat\Sigma)\). Under misspecification or non-normality,
\(\Gamma\) is the actual fourth-moment covariance.

For independent sample groups \(b=1,\ldots,B\), stack
\(\hat\sigma=(\hat\sigma_1\T,\ldots,\hat\sigma_B\T)\T\). With
\(n_b/N\to\pi_b\),
\[
  \sqrt N(\hat\sigma-\sigma_0)
    \Rightarrow
  N\{0,\operatorname{blockdiag}(\Gamma_b/\pi_b)\}.
\]
All formulas below apply to this stacked vector after replacing \(\Gamma\) by
the block diagonal covariance.

\section{The Guttman Map}

Fix a simple CFA loading pattern. Let \(X\in\{0,1\}^{p\times k}\) be the
indicator-by-factor membership matrix, with one nonzero entry per row in the
clean simple-structure case. Let \(S(\sigma)\) be the symmetric matrix whose
lower triangle is \(\sigma\).

Guttman's method first replaces the diagonal of \(S\) by communalities. Write
\[
  h(\sigma)\in\mathbb{R}^p,\qquad
  C(\sigma)=S(\sigma)-\diag\{\diag S(\sigma)\}+\diag\{h(\sigma)\}.
\]
For the Spearman-Harman rule used by lavaan and magmaan starts, the unclamped
correlation-scale communality for indicator \(i\) in its factor block is
\[
  h_i^{\mathrm{cor}}(R)
    =
  {|P_i|}^{-1}\sum_{(j,\ell)\in P_i}
    \frac{R_{ij}R_{i\ell}}{R_{j\ell}},
  \qquad
  h_i(\Sigma)=\Sigma_{ii}h_i^{\mathrm{cor}}\{R(\Sigma)\},
\]
where \(P_i=\{(j,\ell):j<\ell,\ j\ne i,\ \ell\ne i\}\). The production start
routine clamps this value and falls back in singular cases; the smooth theory
below assumes an interior point where no clamp or fallback is active. A
different communality rule only changes \(Dh(\sigma_0)\).

Define
\[
  A=CX,\qquad
  B=X\T C X,\qquad
  D=\diag\{B\},\qquad
  Q=D^{-1/2}.
\]
Then
\[
  L=AQ,\qquad
  P=QBQ,\qquad
  K^\star=LP^{-1}.
\]
If \(m_f\) is the marker indicator for factor \(f\), set
\[
  M=\diag(K^\star_{m_1 1},\ldots,K^\star_{m_k k}).
\]
The marker-scaled Guttman estimates are
\[
  \Lambda_G=K^\star M^{-1},
  \qquad
  \Phi_G=MPM,
  \qquad
  \psi_G=\diag\{S-\Lambda_G\Phi_G\Lambda_G\T\}.
\]
The reported parameter map is a fixed selection
\[
  \tau_G(\sigma)
    =
  \bigl[
    \operatorname{sel}_{\Lambda}(\Lambda_G)\T,\ 
    \vech(\Phi_G)\T,\ 
    \psi_G\T
  \bigr]\T .
\]
This is the version closest to the current \texttt{guttman\_start\_values()}
code: the full \(K^\star\) is computed, but only the model's allowed loading
cells are read out. Dhaene and Rosseel mention a constrained quadratic program
for enforcing simple structure. If that projection is used instead, the same
delta method applies after differentiating the fixed-active-set KKT system.

\section{Guttman Influence Function}

Let \(dS\) be a symmetric covariance perturbation and write all quantities at
\(\sigma_0\). Let
\[
  dC=dS-\diag\{\diag(dS)\}+\diag\{Dh[\sigma_0]\,d\sigma\}.
\]
The differential chain is
\[
\begin{aligned}
  dA &= dC\,X,\\
  dB &= X\T dC\,X,\\
  dD &= \diag\{dB\},\\
  dQ &= -\frac12 D^{-3/2}dD,\\
  dL &= dA\,Q+A\,dQ,\\
  dP &= dQ\,B\,Q+Q\,dB\,Q+Q\,B\,dQ,\\
  dK^\star &= dL\,P^{-1}-L\,P^{-1}dP\,P^{-1}.
\end{aligned}
\]
Next
\[
  dM=\diag(dK^\star_{m_1 1},\ldots,dK^\star_{m_k k}),
\]
and
\[
\begin{aligned}
  d\Lambda_G
    &= dK^\star M^{-1}-K^\star M^{-1}dM M^{-1},\\
  d\Phi_G
    &= dM\,P\,M+M\,dP\,M+M\,P\,dM,\\
  d\psi_G
    &=
  \diag\{dS
    -d\Lambda_G\Phi_G\Lambda_G\T
    -\Lambda_G d\Phi_G\Lambda_G\T
    -\Lambda_G\Phi_G d\Lambda_G\T\}.
\end{aligned}
\]

\begin{proposition}[Guttman delta method]\label{prop:guttman}
Assume \(\Gamma\) is finite, \(h\) is differentiable at \(\sigma_0\), the
diagonal of \(B\) is positive, \(P\) is nonsingular, and all marker entries in
\(M\) are nonzero. Then
\[
  \sqrt n\{\tau_G(\hat\sigma)-\tau_G(\sigma_0)\}
    \Rightarrow
  N(0,\Omega_G),
  \qquad
  \Omega_G=J_G\Gamma J_G\T,
\]
where \(J_G=D\tau_G[\sigma_0]\) is the linear map defined by the differential
chain above. The observation-level influence function is
\[
  \IF_G(y_i)=J_G\,\vech\{(y_i-\mu_0)(y_i-\mu_0)\T-\Sigma_0\}.
\]
\end{proposition}

\begin{remark}
The current magmaan implementation already has most of the numerical map in
\texttt{src/estimate/start\_guttman.cpp}. Turning it into an estimator is not
hard. The open design choice is whether the estimator should be the clean
smooth map above, the lavaan-style clamped start map, or the constrained
projection variant used by Dhaene and Rosseel. For asymptotic comparisons, the
clean map is the right first target.
\end{remark}

\section{Guttman Omega and Alpha}

Let \(w\in\mathbb{R}^p\) be fixed score weights. For a diagonal-residual CFA
model, coefficient omega for the composite \(w\T y\) is
\[
  \omega(\Lambda,\Phi,\psi;w)
    =
  \frac{w\T\Lambda\Phi\Lambda\T w}{w\T S w},
\]
where \(S\) is the covariance matrix whose reliability is being summarized
\cite{mcdonald1999}. Applying this formula to the Guttman map gives
\[
  \omega_G(\sigma;w)
    =
  \frac{w\T\Lambda_G\Phi_G\Lambda_G\T w}{w\T S(\sigma) w}.
\]
This looks parameter-based, but the parameters cancel. Since
\[
  \Lambda_G\Phi_G\Lambda_G\T
    =
  K^\star P (K^\star)\T
    =
  L P^{-1}L\T
    =
  C X(X\T C X)^{-1}X\T C,
\]
we have the covariance-only form
\begin{equation}\label{eq:guttman-omega-general}
  \omega_G(\sigma;w)
    =
  \frac{
    w\T C X(X\T C X)^{-1}X\T C w
  }{
    w\T S(\sigma)w
  }.
\end{equation}
Thus, once the grouping matrix \(X\) and communality rule \(h(\sigma)\) are
fixed, Guttman omega is a population covariance functional. It does not require
iterative loadings, fitted residual variances, or marker choices. It is not
assumption-free: \(X\), \(h\), and the nonsingularity of \(X\T C X\) are part of
the definition.

For one factor, \(X=\mathbf{1}\). Put
\[
  a=C\mathbf{1},\qquad b=\mathbf{1}\T C\mathbf{1},\qquad
  T=\mathbf{1}\T S\mathbf{1}.
\]
Then
\[
  \omega_G(w)=\frac{(w\T a)^2}{b\,w\T S w}.
\]
For the unit-weight total score \(w=\mathbf{1}\), this collapses further:
\begin{equation}\label{eq:guttman-omega-one-factor}
  \omega_G
    =
  \frac{\mathbf{1}\T C\mathbf{1}}{\mathbf{1}\T S\mathbf{1}}
    =
  \frac{\sum_{i\ne j}S_{ij}+\sum_i h_i}{T}
    =
  1-\frac{\sum_i(S_{ii}-h_i)}{T}.
\end{equation}
Here \(h_i\) is the communality placed on item \(i\)'s diagonal in \(C\). For
the Spearman-Harman/Guttman rule above, in the one-factor case,
\[
  \begin{aligned}
  h_i
    &=
  {\binom{p-1}{2}}^{-1}
  \left\{
    \sum_{1\le j<\ell\le p}
      \frac{S_{ij}S_{i\ell}}{S_{j\ell}}
    -(p-1)S_{ii}
  \right\} \\
    &=
  \{2\binom{p-1}{2}\}^{-1}
  \left\{
    \sum_{j,\ell=1}^{p}
      \frac{S_{ij}S_{i\ell}}{S_{j\ell}}
    -\sum_{j=1}^{p}\frac{S_{ij}^2}{S_{jj}}
    -2(p-1)S_{ii}
  \right\}.
  \end{aligned}
\]
This is the covariance-scale form of the correlation expression
\(S_{ii}R_{ij}R_{i\ell}/R_{j\ell}\): the item variance factors cancel, and the
unrestricted pair sum contains \(p-1\) pairs involving item \(i\), each reducing
to \(S_{ii}\). The full double sum counts each unordered off-diagonal pair
twice. Its subtraction removes the diagonal contribution
\(\sum_jS_{ij}^2/S_{jj}\) and the \(2(p-1)\) off-diagonal terms involving item
\(i\), each reducing to \(S_{ii}\). In a multi-factor simple-structure block,
replace \(p\) by the block size \(p_f\) and sum only within the block. Under an
exact congeneric one-factor population, this population rule recovers
\(h_i=\lambda_i^2\phi\).
So the one-factor Guttman-derived omega is especially simple: replace each
item variance by its Guttman communality, sum all entries, and divide by the
observed total-score variance.

Cronbach's alpha for the same unit-weight score is
\cite{cronbach1951}
\begin{equation}\label{eq:alpha}
  \alpha
    =
  \frac{p}{p-1}
  \left(
    1-\frac{\tr(S)}{\mathbf{1}\T S\mathbf{1}}
  \right)
    =
  \frac{p}{p-1}\,
  \frac{\sum_{i\ne j}S_{ij}}{T}.
\end{equation}
Equations~\eqref{eq:guttman-omega-one-factor} and \eqref{eq:alpha} show the
contrast. Alpha extrapolates reliability from the off-diagonal covariance
average. Guttman omega uses the same off-diagonal covariance sum, but adds an
explicit diagonal common-variance estimate \(\sum_i h_i\).

Under an exact one-factor congeneric population,
\[
  \Sigma=\lambda\lambda\T\phi+\diag(\psi),
\]
and under a communality rule that recovers \(h_i=\lambda_i^2\phi\),
\[
  \omega_G
    =
  \frac{\phi(\mathbf{1}\T\lambda)^2}
       {\phi(\mathbf{1}\T\lambda)^2+\mathbf{1}\T\psi},
\]
which is the usual omega total for the unit-weight sum score. In the same
case,
\[
  \omega_G-\alpha
    =
  \frac{
    \phi\{p\sum_i\lambda_i^2-(\sum_i\lambda_i)^2\}
  }{
    (p-1)T
  }
  \ge 0,
\]
with equality exactly when the keyed loadings are equal. Thus alpha and
Guttman omega agree in the tau-equivalent case; with congeneric loadings,
alpha is lower because it has no direct diagonal common-variance estimate.

Away from the exact one-factor model, both quantities remain smooth covariance
functionals, but they target different summaries. Alpha needs no loading
pattern and is easiest to compute. Guttman omega is still non-iterative, but it
imports the CFA grouping pattern and the chosen communality estimator. Its
large-sample variance follows immediately by the delta method:
\[
  \sqrt n\{\hat\omega_G-\omega_G\}
    \Rightarrow
  N\{0,\nabla\omega_G(\sigma_0)\T\Gamma\nabla\omega_G(\sigma_0)\},
\]
and alpha has the same form with \(\nabla\alpha(\sigma_0)\).

\paragraph{Relation to Hancock and An (2020).}
Hancock and An \cite{hancockan2020} independently proposed a closed-form
single-factor omega and evaluated it in a large simulation against the CFA-based
estimate. Their loading uses Spearman's \cite{spearman1927} ratio-of-sums rule,
\[
  \hat\lambda_i
    =
  \left\{
    \frac{\sum_{j<k}S_{ij}S_{ik}}{\sum_{j<k}S_{jk}}
  \right\}^{1/2},
  \qquad
  \hat\omega=\frac{(\sum_i\hat\lambda_i)^2}{T},
  \qquad j,k\ne i.
\]
Under an exact one-factor population this coincides with
\eqref{eq:guttman-omega-one-factor}: both recover \(\hat\lambda_i^2=\lambda_i^2\phi\)
and hence the same omega. Off the one-factor model the two differ. The communality
\(h_i\) above is an \emph{average} of the per-triple ratios
\(S_{ij}S_{i\ell}/S_{j\ell}\), whereas Hancock and An pool numerators and
denominators before dividing; pooling is the more numerically stable choice when
some \(S_{j\ell}\) are near zero, so for the single-factor unit-weight case their
ratio-of-sums aggregation is preferred. The single-factor closed-form omega is
therefore not new with this note (numerically confirmed: the two agree to machine
precision on a one-factor population and diverge, e.g.\ \(0.712\) vs \(0.749\), on a
two-factor covariance analyzed as one factor).

What \eqref{eq:guttman-omega-general} adds beyond Hancock and An, who restrict
to a single factor, is the grouping matrix \(X\) and the weights \(w\): a
closed-form omega \emph{total} for multi-factor, multi-group, and
weighted-composite reliability. The natural next target is omega
\emph{hierarchical}: a second-stage centroid on the Guttman factor correlation
\(\Phi_G\) (a Schmid-Leiman general factor) yields closed-form general-factor
loadings whenever \(k\ge 3\), and hence a closed-form \(\omega_h\); a
non-proportional bifactor solution lies outside the simple-structure \(X\) and is
not available by this route.

\section{Guttman Omega Hierarchical}

Omega hierarchical \(\omega_h\) is the share of total score variance carried by a
single general factor under a Schmid-Leiman \cite{schmidleiman1957}
orthogonalization of a correlated \(k\)-factor model. Write
\(\Sigma=\Lambda\Phi\Lambda\T+\Psi\) with the second-order factorization of the
factor correlation \(\Phi=\gamma\gamma\T+\Psi_\varphi\), where
\(\gamma\in\mathbb{R}^k\) holds the loadings of the first-order factors on the
general factor. With simple structure the general-factor item loadings are
\(\lambda_g=\Lambda\gamma\), and
\[
  \omega_h=\frac{(\mathbf 1\T\lambda_g)^2}{\mathbf 1\T S\mathbf 1}
    =\frac{(\mathbf 1\T\Lambda\gamma)^2}{T},\qquad T=\mathbf 1\T S\mathbf 1 .
\]

The first stage is the Guttman map already given: in the unit-variance metric the
loadings are \(K^\star=CXB^{-1}D^{1/2}\) and the factor correlation is
\(P=D^{-1/2}BD^{-1/2}\), with \(B=X\T C X\), \(D=\operatorname{Diag}B\). The second
stage applies the \emph{same} centroid extraction to \(P\): replace its unit
diagonal by second-order communalities to form \(C_P\), and read off
\[
  \gamma_P=\frac{C_P\mathbf 1}{\sqrt{\mathbf 1\T C_P\mathbf 1}},\qquad
  \lambda_g=K^\star\gamma_P .
\]
Then \(\lambda_g\lambda_g\T=K^\star\gamma_P\gamma_P\T K^{\star\T}\) is the
general-factor item covariance and
\begin{equation}\label{eq:omega-h}
  \omega_h(\sigma)
    =\frac{(\mathbf 1\T K^\star\gamma_P)^2}{T}
    =\frac{\bigl(a\T B^{-1}D^{1/2}C_P\mathbf 1\bigr)^2}
          {(\mathbf 1\T C_P\mathbf 1)\,T},\qquad a=X\T C\mathbf 1 .
\end{equation}
This needs \(k\ge3\) first-order factors, since the second-stage triples are over
factors. It is the Schmid-Leiman \emph{proportional} general factor; a free
(non-proportional) bifactor is outside the simple-structure \(X\).

In Hancock and An's \cite{hancockan2020} elementary form this is their
ratio-of-sums rule applied at two levels. Within block \(f\),
\(\hat\lambda_i=\{\sum_{j<\ell\in f}S_{ij}S_{i\ell}/\sum_{j<\ell\in f}S_{j\ell}\}^{1/2}\)
and \(s_f=\sum_{i\in f}\hat\lambda_i\); the factor correlations are
\(\hat\varphi_{fg}=\sum_{i\in f,\,j\in g}S_{ij}/(s_fs_g)\); the same rule one level
up gives
\(\hat\gamma_f=\{\sum_{g<h\ne f}\hat\varphi_{fg}\hat\varphi_{fh}/\hat\varphi_{gh}\}^{1/2}\);
and \(\hat\omega_h=(\sum_f\hat\gamma_fs_f)^2/T\). On the exact higher-order
manifold this two-level form and \eqref{eq:omega-h} coincide and recover the true
\(\omega_h\).

\paragraph{Influence function.} \(\omega_h=g(\sigma)\) is a smooth covariance
functional, so by the delta method
\(\sqrt n\{\hat\omega_h-\omega_h\}\Rightarrow N(0,\nabla g\T\Gamma\nabla g)\) with
\(\IF(y_i)=\nabla g(\sigma_0)\T z_i\). The gradient composes the matrix
differentials
\[
\begin{aligned}
  dC &= dS-\operatorname{Diag}(\operatorname{diag}dS)+\operatorname{Diag}(dh),\\
  dB &= X\T\,dC\,X,\\
  dP &= D^{-1/2}\,dB\,D^{-1/2}-\tfrac12(\Delta P+P\Delta),
        \quad \Delta=D^{-1}\operatorname{Diag}(dB),\\
  dC_P &= dP-\operatorname{Diag}(\operatorname{diag}dP)+\operatorname{Diag}(dh^\varphi),
\end{aligned}
\]
where \(dh\) and \(dh^\varphi\) are the first- and second-stage communality
differentials, the only stage-specific pieces; a different communality rule
changes nothing else. Vectorizing, each stage is Kronecker-linear in
\(\operatorname{vech}dS\): the lift is \(\operatorname{vec}dB=(X\T\otimes
X\T)\operatorname{vec}dC\), the inverse uses \(d(B^{-1})=-B^{-1}(dB)B^{-1}\), and
every scalar contraction uses \(a\T(dM)b=(b\T\otimes a\T)\operatorname{vec}dM\).
With \(N=a\T B^{-1}D^{1/2}C_P\mathbf 1\),
\[
\begin{aligned}
  dN ={}& (X\T dC\,\mathbf 1)\T B^{-1}D^{1/2}C_P\mathbf 1
     - a\T B^{-1}(dB)B^{-1}D^{1/2}C_P\mathbf 1 \\
     &+ a\T B^{-1}(dD^{1/2})C_P\mathbf 1
     + a\T B^{-1}D^{1/2}(dC_P)\mathbf 1,
\end{aligned}
\]
and the quotient rule on \(\omega_h=N^2/[(\mathbf 1\T C_P\mathbf 1)T]\) gives
\(\nabla\omega_h\). Only the lift \(X\T\otimes X\T\) and the two communality blocks
distinguish this from the single-factor case, so \(\Gamma\)'s vech structure
passes straight into \(\Omega=\nabla\omega_h\T\Gamma\nabla\omega_h\). Numerically,
\eqref{eq:omega-h} recovers the true \(\omega_h\) to machine precision on a
higher-order population (\(k=3\) blocks of three items, \(\omega_h=0.6134277\)),
and its delta-method standard error matches a Monte Carlo standard deviation to
within one percent.

\section{Moment-Quadratic Estimators}

Let \(m(\theta)\) be the model-implied covariance moment vector and
\[
  r(\theta,\sigma)=m(\theta)-\sigma.
\]
A fixed-weight ULS/GLS/WLS estimator minimizes
\[
  q_W(\theta,\sigma)=\frac12 r(\theta,\sigma)\T W r(\theta,\sigma),
\]
with pseudo-true value \(\theta_W\). Let
\[
  \Delta=\frac{\partial m(\theta_W)}{\partial\theta\T},
  \qquad
  H_W
    =
  \Delta\T W\Delta
    +\sum_{\ell}\{Wr(\theta_W,\sigma_0)\}_{\ell}
       \frac{\partial^2 m_\ell(\theta_W)}
            {\partial\theta\,\partial\theta\T}.
\]
The second term is the misspecification curvature term. It vanishes at exact
fit and is easy to miss if one silently substitutes Fisher or Gauss-Newton
information.

For fixed \(W\),
\[
  \IF_W(y_i)=
  H_W^{-1}\Delta\T W z_i,
  \qquad
  \Omega_W
    =
  H_W^{-1}\Delta\T W\Gamma W\Delta H_W^{-1}.
\]
If \(W=W(\hat\sigma)\), the weight influence is first-order under
misspecification:
\[
  \IF_{W(\hat\sigma)}(y_i)
    =
  H_W^{-1}\Delta\T W z_i
  -
  H_W^{-1}\Delta\T \{DW[\sigma_0]z_i\}\,r(\theta_W,\sigma_0).
\]
This second term drops out at exact fit. It is therefore most important in the
setting the user asked about: fixed model misspecification with GLS/WLS style
weights. This is also the term already handled by magmaan's continuous LS
influence machinery when the caller asks for data-dependent weight influence.

The efficient covariance-moment estimator for a fixed covariance target uses
\(W=\Gamma^{-1}\), giving the familiar
\[
  \Omega_{\mathrm{eff}}
    =
  (\Delta\T\Gamma^{-1}\Delta)^{-1}
\]
only at exact fit or when the curvature and weight-influence terms are absent.
Under misspecification the observed bread, not the Fisher simplification, is the
general object.

\section{Normal-Theory ML}

For covariance-only NT estimation,
\[
  F_{\mathrm{NT}}(\theta,\sigma)
    =
  \log|\Sigma_\theta|
  +\tr\{S(\sigma)\Sigma_\theta^{-1}\}
  -\log|S(\sigma)|-p.
\]
The term \(-\log|S|\) is irrelevant for the first-order condition. At the
pseudo-true point \(\theta_{\mathrm{NT}}\), define
\[
  U_a(\theta,\sigma)
    =
  \tr\left[
    \{\Sigma_\theta^{-1}
      -\Sigma_\theta^{-1}S(\sigma)\Sigma_\theta^{-1}\}
    \frac{\partial\Sigma_\theta}{\partial\theta_a}
  \right].
\]
Let \(H_{\mathrm{NT}}=\partial U/\partial\theta\T\) be the observed
pseudo-true bread. For a covariance perturbation \(dS\),
\[
  \{G_{\mathrm{NT}}d\sigma\}_a
    =
  -\tr\left[
    \Sigma_\theta^{-1}dS\Sigma_\theta^{-1}
    \frac{\partial\Sigma_\theta}{\partial\theta_a}
  \right].
\]
Thus
\[
  \IF_{\mathrm{NT}}(y_i)
    =
  -H_{\mathrm{NT}}^{-1}G_{\mathrm{NT}}z_i.
\]
When the normal model is correct this is the usual efficient likelihood
influence function. Under non-normality or fixed structural misspecification,
the sandwich \(H^{-1}G\Gamma G\T H^{-1}\) is the robust covariance around the
NT pseudo-true parameter.

\section{What Efficiency Means Under Misspecification}

Under fixed misspecification, Guttman, NT, ULS, GLS, and WLS generally estimate
different functionals of the same population covariance. A variance ratio is
therefore not automatically an MSE ratio.

For a scalar contrast \(a_e\{\tau_e(P)\}\), the variance-only ratio is
\[
  \rho_{G/e}
    =
  \frac{\nabla a_G\T \Omega_G \nabla a_G}
       {\nabla a_e\T \Omega_e \nabla a_e}.
\]
This is meaningful when the contrast names the same scientific quantity and
the pseudo-targets are close enough that first-order variance is the object of
interest.

For fixed misspecification, a finite-\(n\) target-risk comparison is more
honest:
\[
  R_e(n;A,\tau_0)
    =
  \{\tau_e(P)-\tau_0\}\T A\{\tau_e(P)-\tau_0\}
  +\frac1n\tr(A\Omega_e).
\]
Dhaene and Rosseel's MSE comparisons are of this type, with \(A\) selecting
groups of parameters and \(\tau_0\) the data-generating parameter vector. An
analytic experiment should report both \(R_e\) and the variance-only ratios.

A third option is local misspecification:
\[
  \sigma_n=\sigma\{\theta_0\}+n^{-1/2}\delta.
\]
Then bias and variance both contribute at first order, and the influence maps
above give closed-form local asymptotic risk. This is probably the cleanest
paper-level efficiency comparison.

\section{Computational Accounting}

The robust SEs for ordinary NT/ULS estimators and for the Guttman covariance
map use different Jacobians. NT/ULS inference starts from the model moment
Jacobian
\[
  \Delta(\theta)=\frac{\partial\sigma(\theta)}{\partial\theta\T}.
\]
Once \(\Delta\) is built, the bread and residual projector are low-dimensional
linear algebra in the parameter and residual subspaces. In contrast, Guttman
map-delta SEs require
\[
  J(\sigma)=\frac{\partial\tau(\sigma)}{\partial\vech(S)\T},
\]
because the estimator itself is a direct covariance functional. This is a
different object, but it is not a fundamental reason for SEs to be orders of
magnitude slower than NT/ULS.

The first implementation bundled three responsibilities in one call:
standard errors, the residual goodness-of-fit statistic, and the reference
weighted-\(\chi^2\) spectrum. The spectrum forms
\[
  M = I-\Delta J,\qquad U=M\T V M,
\]
on the full \(p^\star\times p^\star\) moment space and then diagonalizes the
profile contrast \(U\Gamma\). At \(p=25\), \(p^\star=325\), so this dense
spectral step dominates a call that only reads the SEs. The same bundle also
materialized empirical \(\Gamma\) as a dense \(p^\star\times p^\star\) matrix
before contracting \(J\Gamma J\T\), whereas an SE-only path can stream
casewise moment rows
\[
  u_i = J\{\vech[(x_i-\bar x)(x_i-\bar x)\T]-\vech(S)\}
\]
and accumulate \(\sum_i u_i u_i\T/N^2\) directly in parameter space.

The implementation now separates these responsibilities. The
\texttt{noniterative\_se} family computes only \(J\Gamma J\T/N\) and skips the
\(I-\Delta J\) residual projector and weighted-\(\chi^2\) spectrum. Its
empirical overloads use the casewise contraction above, so they never allocate
dense \(\Gamma\). The full \texttt{noniterative\_inference} family still builds
the GOF bundle because the scaled and mixture reference tests genuinely need the
moment-space spectrum.

\paragraph{Timing note, 2026-07-09.}
After the SE-only split and the first batched configural Jacobian, the remaining
runtime at \(p=25\) was not the empirical meat. In the experiment-58 normal
shape with five factors, five indicators per factor, and empirical SEs, the R
harness reported configural Guttman SE medians of about \(5\) ms for
\texttt{guttman\_std}, \texttt{guttman\_unit}, and \texttt{guttman\_gls}. A
repeated-call microbenchmark on the same shape gave roughly \(4.6\) ms for
\texttt{standardized}/\texttt{unit} and \(5.3\) ms for
\texttt{gls\_aligned}, with normal-theory SEs essentially the same cost. Varying
\(n\) from \(30\) to \(3000\) showed that empirical contraction only becomes
visible for large \(n\). The fixed cost was therefore the estimator-map
Jacobian.

Profiling the Jacobian itself found two avoidable costs. First, the block-GMM
communality Jacobian formed derivative rows at the full global covariance width
and then kept only the block-active columns; at \(p=25\), each five-indicator
block needs \(15\) columns, not the global \(325\). Building
\(d\Gamma_{\mathrm{cor}}\) directly on the active columns removed most of the
communality-Jacobian time. Second, the GMM derivative materialized every dense
\(dW\), even though
\[
  d(A\T Wb)-d(A\T W A)h
    =
  dA\T W e + A\T dW e + A\T W(db-dA\,h),
  \qquad e=b-Ah,
\]
so the weight derivative is needed only through \(dW e\). Using the fixed-rank
pseudo-inverse derivative,
\[
  dW e =
  -W\,d\Omega\,We
  +W^2 d\Omega (I-\Omega W)e
  +(I-W\Omega)d\Omega\,W^2e,
\]
and applying \(d\Omega x=dM\,\Gamma M\T x+M\,d\Gamma M\T x+
M\,\Gamma dM\T x\) avoids constructing \(dW\) matrices. A guarded
full-column-rank fast path for \((M\Gamma M\T)^+\) gives smaller gains by using
\[
  (M\Gamma M\T)^+ =
  M(M\T M)^{-1}\Gamma^{-1}(M\T M)^{-1}M\T
\]
when \(M\) and \(\Gamma\) are full rank, with the eigensolve pseudo-inverse as
fallback. After these changes the same \(p=25\) repeated-call probe reports
about \(1.5\) ms for \texttt{standardized}/\texttt{unit} and \(2.3\) ms for
\texttt{gls\_aligned}. The remaining cost is no longer dominated by the
communality GMM derivative; it is split between the communality Jacobian and the
downstream score-regression derivative, with the GLS-aligned composite-weight
derivative accounting for the extra time.

For a covariance block with \(p^\star=p(p+1)/2\), the old configural analytic
implementation computed \(p^\star\) directional derivatives. Each direction
rebuilt the correlation-scale communality derivative, the triad-GMM weight
derivative, composite-weight derivative, score covariance inverse, marker
scaling derivative, and residual split. The point estimator itself was cheap;
the SE bottleneck was repeating this derivative work over every covariance
coordinate.

For the GLS-aligned configural Guttman map the regular-interior derivative can
be batched. The blockwise GMM communality solves
\[
  h=(A\T W A)^{-1}A\T Wb,\qquad
  dh=(A\T W A)^{-1}\{d(A\T Wb)-d(A\T W A)h\},
\]
or \(dh=A^{-1}(db-dA\,h)\) in just-identified blocks. Computing all columns of
\(Dh/D\vech(S)\) at once lets the implementation reuse the base correlation
rows, \(A,b,W\), the normal-equation factorization, the fixed-rank
Moore--Penrose derivative of \(W\), and the downstream
\(H,B,Q,K^\star,\Lambda,\Phi,\psi\) matrices.

Restricted Guttman is now analytic on the same regular-interior principle.
Active residual rows enter the communality problem as \(Rh=r\). With
\[
  G=A\T W A,\qquad c=A\T Wb,
\]
the full-rank constrained derivative solves
\[
  \begin{bmatrix}G&R\T\\R&0\end{bmatrix}
  \begin{bmatrix}dh\\d\lambda\end{bmatrix}
  =
  \begin{bmatrix}
    dc-dG\,h-dR\T\lambda\\
    dr-dR\,h
  \end{bmatrix}.
\]
Loading equality rows are differentiated through the same composite-metric
projection used by the point estimator. Rank-deficient pseudo-inverse points,
singular loading-projection metrics, and other boundary cases fall back to the
older central finite-difference restricted map. Thus the regular restricted
path avoids \(2p^\star\) full refits, while preserving the conservative
fallback needed for active-set or rank changes.

The first restricted implementation still assembled the regular derivative one
covariance coordinate at a time. On the same \(p=25\) shape this made a
residual-restricted SE call about \(166\) ms, even though no finite-difference
refits were being used. The default \texttt{gmm\_block} restricted path first
batched the constrained \(h^2\) derivative: it builds the block-active GMM RHS
columns with the same \(dW e\) action used by the configural path, solves the
KKT system once with all right-hand sides, and then feeds the constrained
\(dH\) diagonal into the batched score-regression derivative. Loading-only
restrictions also batch the composite-metric projection derivative, reusing the
base \(W^{-1}\), \((RW^{-1}R\T)^{-1}\), and projection matrix. A repeated-call
\(p=25\) probe reports roughly \(1.8\) ms for residual-only restrictions,
\(2.8\) ms for loading-only restrictions, and \(3.2\) ms for both under
\texttt{unit}/\texttt{standardized}; the corresponding \texttt{gls\_aligned}
numbers are about \(2.7\), \(3.7\), and \(3.9\) ms.

The restricted batching now covers the rest of the regular LS-form restricted
path except \texttt{gmm\_full}. For \texttt{triad\_ls} and
\texttt{anchor\_triad\_ls}, the full normal-equation RHS
\(dc-dG\,h\) is accumulated directly from all local covariance columns; for
\texttt{gmm\_block}, the same RHS is built with block-active GMM weight actions.
Grouped restricted maps embed the active sample block's RHS in the stacked
communality system, solve the joint constrained KKT system once, and pass the
resulting cross-block \(dH\) diagonals to every block's batched score-regression
derivative. This is the same chain rule as the earlier directional grouped
implementation, but without rebuilding the communality system for each
covariance coordinate. \texttt{gmm\_full} still uses the direction-wise analytic
fallback; making it competitive requires a batched derivative of the joint
selected-triad weight.

A later restricted-cost pass removed a smaller remaining duplication. The
loading projection applies
\[
  dl = (I-CR)\,dl_0 + (I-CR)\,W^{-1}(dW\,a),\qquad
  C=W^{-1}R\T(RW^{-1}R\T)^{-1},\quad a=C(Rl_0-r),
\]
so no per-column \(dW^{-1}\), \(d(RW^{-1}R\T)^{-1}\), or \(dC\) matrices are
materialized. In a p = 25 restricted empirical-SE probe, absolute timings
depended on the build tree, but \texttt{unit}/\texttt{standardized} were about
\(1.5\times\) the ULS robust-SE time and \texttt{gls\_aligned} about
\(1.7\times\), meeting the current ``within \(2\times\) NT'' target.

\section{Proposed magmaan Experiment}

The experiment should be deterministic and population based first, with Monte
Carlo used only to validate the formulas.

\paragraph{1. Estimator maps.}
Implement a frontier Guttman CFA estimator that returns
\((\Lambda_G,\Phi_G,\psi_G)\) per covariance block using the clean map in this
note. Reuse \texttt{cfa\_block\_layouts()} and
\texttt{guttman\_start\_values()} as implementation scaffolding, but keep
finite-sample clamps out of the analytic target.

\paragraph{2. Jacobians.}
Start with central finite differences of \(\tau_G(\sigma)\) to verify the
chain above. Then add the analytic differential for speed and numerical
diagnostics. The Spearman-Harman communality derivative is a quotient-rule
calculation through the covariance-to-correlation map.

\paragraph{3. Comparator estimators.}
For each population covariance, fit the analysis model at the population
moments with NT, ULS, GLS, DWLS/WLS, and possibly empirical-ADF WLS. Use the
current observed-bread sandwich machinery for NT and continuous LS. For
data-dependent GLS/WLS weights, run both fixed-weight and weight-influence
versions to show how much the misspecification term matters.

\paragraph{4. Designs.}
Replicate the Dhaene-Rosseel three-factor, three-indicator design with
reliability \(0.5/0.8\), factor correlations \(0.1/0.4\), omitted
cross-loadings, and omitted residual covariances. Add one larger block design
with six to twelve indicators per factor, because Guttman's computational
appeal grows with \(p\).

\paragraph{5. Distributional covariance.}
Evaluate \(\Gamma\) under three choices: normal theory, an analytic elliptical
or independent-component heavy-tailed fourth-moment model, and empirical
fourth moments from a very large simulated DGP draw. This separates structural
misspecification from non-normal covariance noise.

\paragraph{6. Outputs.}
For each parameter group and selected scalar contrasts, report:
\[
  \diag(\Omega_G)/\diag(\Omega_e),\qquad
  \tr(A\Omega_G)/\tr(A\Omega_e),\qquad
  R_G(n;A,\tau_0)/R_e(n;A,\tau_0).
\]
Also report the pseudo-target displacement
\(\tau_e(P)-\tau_0\), because under fixed misspecification this often dominates
large-\(n\) MSE.

\paragraph{7. Validation.}
For a small subset, simulate many datasets, compute all estimators, and check
that \(n\,\widehat{\Cov}(\hat\tau_e)\) matches \(\Omega_e\). For Guttman this
is the main new validation. For NT and LS it should agree with existing
misspecification-robust sandwich code. This is implemented as the
\texttt{--validate} path in experiment 40: at the well-specified cell with
\(2000\) replications and \(N=1000\), the Guttman all-parameter asymptotic
SD matches simulation to within \(1\%\) (\(-0.8\%\)), bracketed by the
known-correct NT (\(-0.7\%\)) and oracle-WLS (\(-1.0\%\)) sandwiches, so the
finite-difference Jacobian of the Guttman map is as accurate as those
comparators.

\section{Takeaways}

Guttman's estimator has straightforward regular asymptotics because it is a
smooth covariance-summary map away from clamps, singular factor-score
correlations, and zero marker loadings. Its influence function is simply the
Jacobian of that map times the covariance influence row.

The comparison to NT, ULS, GLS, and WLS should not be framed as one universal
efficiency ranking. Under correct normal specification NT/GLS have the usual
optimality story. Under fixed misspecification, the relevant objects are
pseudo-targets, observed breads, empirical meats, and weight influence. That is
exactly the regime where an analytic magmaan experiment can be more informative
than another Monte Carlo table.

\begin{thebibliography}{9}

\bibitem{bollen1989}
Bollen, K. A. (1989).
\newblock \emph{Structural equations with latent variables}.
\newblock Wiley.

\bibitem{browne1974}
Browne, M. W. (1974).
\newblock Generalized least squares estimators in the analysis of covariance
structures.
\newblock \emph{South African Statistical Journal}, 8, 1--24.

\bibitem{cronbach1951}
Cronbach, L. J. (1951).
\newblock Coefficient alpha and the internal structure of tests.
\newblock \emph{Psychometrika}, 16, 297--334.
\newblock \href{https://doi.org/10.1007/BF02310555}
{doi:10.1007/BF02310555}.

\bibitem{dhaene2024}
Dhaene, S., \& Rosseel, Y. (2024).
\newblock An evaluation of non-iterative estimators in confirmatory factor
analysis.
\newblock \emph{Structural Equation Modeling: A Multidisciplinary Journal},
31(1), 1--13.
\newblock \href{https://doi.org/10.1080/10705511.2023.2187285}
{doi:10.1080/10705511.2023.2187285}.

\bibitem{guttman1952}
Guttman, L. (1952).
\newblock Multiple group methods for common-factor analysis: Their basis,
computation, and interpretation.
\newblock \emph{Psychometrika}, 17, 209--222.
\newblock \href{https://doi.org/10.1007/BF02288783}
{doi:10.1007/BF02288783}.

\bibitem{hancockan2020}
Hancock, G. R., \& An, J. (2020).
\newblock A closed-form alternative for estimating \(\omega\) reliability under
unidimensionality.
\newblock \emph{Measurement: Interdisciplinary Research and Perspectives}, 18(1),
1--14.
\newblock \href{https://doi.org/10.1080/15366367.2019.1656049}
{doi:10.1080/15366367.2019.1656049}.

\bibitem{harman1976}
Harman, H. H. (1976).
\newblock \emph{Modern factor analysis} (3rd ed.).
\newblock University of Chicago Press.

\bibitem{spearman1927}
Spearman, C. (1927).
\newblock \emph{The abilities of man: Their nature and measurement}.
\newblock Macmillan.

\bibitem{mcdonald1999}
McDonald, R. P. (1999).
\newblock \emph{Test theory: A unified treatment}.
\newblock Lawrence Erlbaum.

\bibitem{schmidleiman1957}
Schmid, J., \& Leiman, J. M. (1957).
\newblock The development of hierarchical factor solutions.
\newblock \emph{Psychometrika}, 22(1), 53--61.
\newblock \href{https://doi.org/10.1007/BF02289209}
{doi:10.1007/BF02289209}.

\end{thebibliography}

\end{document}
